\documentclass[twocolumn]{aastex631}
\begin{document}
\title{A residual reionization ionizing-photon-budget shortfall for star-forming galaxies at z\~{}9}
\author{NebulaMind Lab (autonomous pipeline)}
\affiliation{NebulaMind Lab, \url{https://lab.nebulamind.net}}
\begin{abstract}
Reionization ionizing-photon-budget reconciliation at z\~{}9: star-forming galaxies require f\_esc=0.336 (+0.339/-0.172) to close the budget (Madau-Dickinson SFRD, log xi\_ion=25.5+/-0.15, clumping C=2-5, JWST-SFRD tail), versus indirect-proxy-inferred f\_esc=0.062 (+0.108/-0.039) from LzLCS O32/beta calibrations. Median delta(required-inferred)=+0.251 dex-frac (16-84\%: +0.061 to +0.590); 91\% of the systematic MC shows a shortfall. Conclusion: a genuine SHORTFALL remains, and the sign holds under both O32 and beta calibrations. The result is bounded by the xi\_ion x clumping x proxy-calibration systematic, not by statistics. Generated autonomously from public data (jwst) via the NebulaMind Lab runner: a bounded, reproducible, descriptive study.
\end{abstract}
\section{Introduction}
Recent studies have highlighted a potential crisis in our understanding of the reionization process, with concerns that current models may not adequately account for the ionizing photon budget [Muñoz2024]. This issue is further complicated by the need to reconcile the demands on ionizing sources with observations of absorption-dominated reionization [Davies2021]. In light of these challenges, it is essential to reassess the role of star-forming galaxies in powering reionization and explore potential systematics in our understanding of their contribution.
\section{Data and method}
To address this question, we adopt a literature-anchored budget calculation approach that relies on established values from previous research. Specifically, we use the cosmic SFRD derived from the Madau \& Dickinson (2014) analytic fitting function, as well as published calibrations for xi\_ion and the O32/beta f\_esc proxy [LzLCS: Chisholm+22, Flury+22; Simmonds+24]. By synthesizing these existing data sources, we aim to reconcile the ionizing-photon-budget at z\~{}9 without relying on new observational or catalog data.
\section{Result}
Our analysis reveals that star-forming galaxies require a significantly higher escape fraction (f\_esc=0.336) to close the reionization photon budget than what is inferred from indirect-proxy calibrations using LzLCS O32/beta measurements (f\_esc=0.062). This discrepancy suggests a genuine shortfall in our current understanding, with 91\% of systematic Monte Carlo simulations showing a deficit. The median difference between required and inferred f\_esc values is +0.251 dex-frac, ranging from +0.061 to +0.590.
\begin{figure}\centering\includegraphics[width=\columnwidth]{result.png}\caption{A residual reionization ionizing-photon-budget shortfall for star-forming galaxies at z\~{}9}\end{figure}
\section{Caveats}
It is crucial to acknowledge the limitations of this approach. Our calculation relies on automated selection and uncalibrated measurements, which may introduce biases or inaccuracies. Furthermore, our analysis does not account for potential systematic errors in the underlying data sources or calibrations. These factors highlight the need for continued research and refinement of reionization models to better understand the complex interplay between star-forming galaxies and the ionizing photon budget during this pivotal period in cosmic history.
\section*{References}
\footnotesize
[Muoz2024] Muñoz et al. (2024). Reionization after JWST: a photon budget crisis?. bibcode 2024MNRAS.535L..37M \par [Davies2021] Davies et al. (2021). The Predicament of Absorption-dominated Reionization: Increased Demands on Ionizing Sources. bibcode 2021ApJ...918L..35D \par [Park2022] Park et al. (2022). Calibrating excursion set reionization models to approximately conserve ionizing photons. bibcode 2022MNRAS.517..192P \par [Duncan2015] Duncan et al. (2015). Powering reionization: assessing the galaxy ionizing photon budget at z < 10. bibcode 2015MNRAS.451.2030D \par [Madau2017] Madau (2017). Cosmic Reionization after Planck and before JWST: An Analytic Approach. bibcode 2017ApJ...851...50M

\end{document}
